\documentclass[11pt,a4paper]{article}
\usepackage[margin=27mm]{geometry}
\usepackage{amsmath,amssymb,amsthm,mathtools,enumitem,indentfirst}
\usepackage[T1]{fontenc}
\usepackage{lmodern}
\usepackage{xcolor}
\usepackage[colorlinks=true,linkcolor=blue,citecolor=blue,urlcolor=blue]{hyperref}
\urlstyle{same}
\setlength{\parindent}{1.5em}
\setlength{\emergencystretch}{2em}
\numberwithin{equation}{section}
\newtheorem{theorem}{Theorem}[section]
\newtheorem{lemma}[theorem]{Lemma}
\newtheorem{proposition}[theorem]{Proposition}
\newtheorem{corollary}[theorem]{Corollary}
\theoremstyle{definition}
\newtheorem{definition}[theorem]{Definition}
\DeclareMathOperator{\Hom}{Hom}
\DeclareMathOperator{\End}{End}
\DeclareMathOperator{\Ext}{Ext}
\DeclareMathOperator{\Aut}{Aut}
\DeclareMathOperator{\Iso}{Iso}
\DeclareMathOperator{\rad}{rad}
\DeclareMathOperator{\soc}{soc}
\DeclareMathOperator{\topm}{top}
\DeclareMathOperator{\add}{add}
\DeclareMathOperator{\im}{im}
\DeclareMathOperator{\coker}{coker}
\DeclareMathOperator{\Tr}{Tr}
\DeclareMathOperator{\gr}{gr}
\newcommand{\FF}{\mathbb F}
\newcommand{\ZZ}{\mathbb Z}
\newcommand{\QQ}{\mathbb Q}
\newcommand{\NN}{\mathbb N}
\newcommand{\R}{\mathcal R}
\newcommand{\gauss}[3]{\genfrac{[}{]}{0pt}{}{#1}{#2}_{#3}}
\title{Hall polynomials for geometrically standard algebras}
\author{Qiyue Tang}
\date{}
\begin{document}
\maketitle
\begingroup\renewcommand{\thefootnote}{}
\footnotetext{2020 Mathematics Subject Classification. Primary 16G60; Secondary 16G20, 16G70.\newline Key words and phrases. Hall polynomial, representation-finite algebra, Auslander defect, mesh category, finite field.}
\endgroup
\begin{abstract}
We prove the existence of integral Hall polynomials for geometrically standard representation-finite algebras over finite fields. In particular, Hall polynomials exist for every representation-finite algebra over a finite field of odd characteristic. The proof expresses Hall numbers through defect modules with prescribed top and socle and computes their automorphism-weighted masses on modulated mesh categories. A Fourier calculation reduces these masses to total masses on periodic Dynkin components; square boundaries are counted by Hall flags of nilpotent semilinear cycles.
\end{abstract}

\section{Introduction}
Let $k=\FF_q$ and let $A$ be a finite-dimensional $k$-algebra of finite representation type. For finite-dimensional right $A$-modules $X,Y,L$ and a finite extension $E/k$, put $M_E=M\otimes_kE$ and define
\[
 F^L_{X,Y}(E)=\#\{U\subseteq L_E:U\simeq Y_E,\ L_E/U\simeq X_E\}.
\]
The first lower index is the quotient and the second is the submodule. A finite extension is \emph{conservative for $A$} if every indecomposable $A$-module remains indecomposable after scalar extension. We say that $A$ has Hall polynomials if, for every $X,Y,L$, there is a polynomial $\varphi^L_{X,Y}\in\ZZ[T]$ such that
\begin{equation}\label{eq:hall}
 F^L_{X,Y}(E)=\varphi^L_{X,Y}(|E|)
 \qquad\text{for every conservative }E/k.
\end{equation}
Hall polynomials for representation-finite hereditary algebras were established by Ringel \cite{Ringel}. Hall polynomials for nilpotent cyclic-quiver representations provide another basic counting input; see \cite[Theorem 16]{Hubery}.

A representation-finite algebra over an algebraically closed field is \emph{mesh-standard} if its category of indecomposable modules is equivalent, with its object labels, to the mesh category of its Auslander--Reiten quiver. We call $A$ \emph{geometrically standard} if $A_{\bar k}$ is mesh-standard. Mesh categories and their finite-field modulation are described in Section~\ref{sec:structure}.

\begin{theorem}\label{thm:main}
Every geometrically standard representation-finite algebra over a finite field has Hall polynomials.
\end{theorem}

\begin{corollary}\label{cor:odd}
Every representation-finite algebra over a finite field of characteristic different from two has Hall polynomials.
\end{corollary}

The structural implication used in this corollary follows from the standard-form theorems reviewed in \cite{Bongartz13}; its passage from the projective points to the complete mesh category is given in Section~\ref{sec:structure}.

The counting argument uses the contravariant defect of a short exact sequence. Its dimension vector determines the alternating sum of the three terms; its top and socle determine the minimal quotient and submodule. This yields an explicit orbit formula for a fixed Hall number. The remaining calculation concerns modules over the stable Auslander algebra, which is a finite modulated partial mesh algebra. We complete its translation intervals to finite cycles, decompose the local linear data into strings and bands, and integrate the quadratic boundary terms. The pure-band contribution is recovered from independently computed total masses on the original periodic components.

\paragraph{Use of generative AI.}
ChatGPT and Codex were used to develop mathematical arguments, perform computational checks, and draft and revise the manuscript.

\section{Defects and Hall numbers}\label{sec:defects}
Fix representatives $I_1,\ldots,I_m$ of the indecomposable $A$-modules, and write
\[
 K_i=\End_A(I_i)/\rad\End_A(I_i)=\FF_{q^{d_i}},\qquad
 D_A=\operatorname{lcm}(d_1,\ldots,d_m).
\]
Finite division rings are fields. Since $k$ is perfect, radicals commute with separable scalar extension. Thus $E=\FF_{q^r}$ is conservative precisely when $\gcd(r,D_A)=1$. Moreover, the modules $(I_i)_E$ form a complete, nonrepeating indecomposable list over every such $E$. Indeed, an $A_E$-module is a direct summand of the scalar extension of its restriction to $A$, because $E/k$ is separable. The restriction is in $\add(\bigoplus_i I_i)$. The semisimple quotient of the endomorphism algebra of this sum becomes $\prod_i(K_i\otimes_kE)$, a product of fields; this also proves that the extended indecomposables remain pairwise nonisomorphic.

Write $h(M,N)=\dim_k\Hom_A(M,N)$ and $e(M)=h(M,M)$. For $M=\bigoplus_i I_i^{n_i}$, set
\begin{equation}\label{eq:aut}
 a_M(T)=T^{e(M)-\sum_i d_i n_i^2}
 \prod_i\prod_{j=0}^{n_i-1}(T^{d_i n_i}-T^{d_i j}).
\end{equation}
Empty products equal one, so $a_0=1$. The formula for the units of a finite algebra, applied to $\End_A(M)$, gives $a_M(|E|)=|\Aut(M_E)|$. All these polynomials are monic.

We use the subring
\[
 \R=\ZZ[T,T^{-1},(T^j-1)^{-1}:j\ge1]\subseteq\QQ(T).
\]
The inverses of the polynomials in \eqref{eq:aut} belong to $\R$.

\begin{lemma}\label{lem:integer}
Let $R\in\R$. If $R(Q)$ is an integer for an unbounded set of positive integers $Q$, at which it is defined, then $R\in\ZZ[T]$.
\end{lemma}
\begin{proof}
Write $R=P/B$ with $P,B\in\ZZ[T]$ and $B$ monic. Division in $\ZZ[T]$ gives $P=BH+U$ with $\deg U<\deg B$. Along the given integers, $U(Q)/B(Q)$ is an integer tending to zero. Hence $U$ has infinitely many zeros and vanishes.
\end{proof}

\subsection{Minimal sequences}
Put $G=\bigoplus_iI_i$ and let
\[
 \Lambda=\End_A(G)/[\operatorname{proj}A],
\]
where $[\operatorname{proj}A]$ is the ideal of maps factoring through projective modules. We use contravariant additive functors on $\add G$. Evaluation at $G$, with precomposition as the right action, identifies the finite-dimensional functors vanishing on projectives with right $\Lambda$-modules. Put $P_M=\Hom_A(-,M)|_{\add G}$.

For a short exact sequence
\[
 \sigma:\quad 0\longrightarrow Y\xrightarrow{i}L\xrightarrow{p}X\longrightarrow0
\]
define $D_\sigma=\coker(P_L\to P_X)$. It has the projective resolution
\begin{equation}\label{eq:resolution}
 0\longrightarrow P_Y\longrightarrow P_L\longrightarrow P_X
 \longrightarrow D_\sigma\longrightarrow0.
\end{equation}
Every finite-dimensional right $\Lambda$-module arises in this way. To see this, take a projective presentation $P_L\to P_X\to D\to0$. Yoneda identifies the first map with a morphism $p:L\to X$. Evaluating at $A$ and using $D(A)=0$ shows that $p$ is surjective. Its kernel supplies \eqref{eq:resolution}.

Every finite projective resolution is the direct sum of a minimal resolution and contractible complexes. Here this follows by eliminating an invertible matrix component of a differential until all differentials lie in the radical of $\add G$. Uniqueness follows successively from uniqueness of projective covers. Consequently, there is a unique minimal sequence
\[
 \sigma_D:\quad 0\to Y_D\to L_D\to X_D\to0
\]
with defect $D$, and every sequence with that defect is uniquely, up to isomorphism, of the form
\begin{equation}\label{eq:contracts}
 \sigma_D\oplus(0\to B\xrightarrow{1}B\to0\to0)
 \oplus(0\to0\to C\xrightarrow{1}C\to0).
\end{equation}
Thus its terms are $Y_D\oplus B$, $L_D\oplus B\oplus C$, and $X_D\oplus C$.

For a nonprojective $I_i$ let $S_i=P_{I_i}/\rad(-,I_i)$. For a defect $D$, write
\[
 \delta_i(D)=\dim_kD(I_i),\quad
 t_i(D)=[\topm D:S_i],\quad s_i(D)=[\soc D:S_i].
\]
The brackets denote multiplicity in a semisimple functor. At projective labels all three numbers are zero. After extension to $E$, the dimension is taken over $E$ and the labels are $(I_i)_E$ and $(S_i)_E$. In particular, $\delta_i$ is divisible by $d_i$. Define
\[
 X(t)=\bigoplus_{I_i\text{ nonprojective}}I_i^{t_i},\qquad
 Y(s)=\bigoplus_{I_i\text{ nonprojective}}(\tau I_i)^{s_i}.
\]
Here $\tau$ denotes the Auslander--Reiten translation from the projective stable category to the injective stable category, and its values are taken without injective summands.

\begin{lemma}\label{lem:ends}
For every defect $D$, one has $X_D\simeq X(t(D))$ and $Y_D\simeq Y(s(D))$.
\end{lemma}
\begin{proof}
The epimorphism $P_{X_D}\to D$ in the minimal resolution is a projective cover, proving the first assertion. Define the covariant defect
\[
 D'_\sigma=\coker\bigl(\Hom_A(L,-)\to\Hom_A(Y,-)\bigr).
\]
Auslander's defect formula gives natural isomorphisms
$D'_\sigma(\tau I)\simeq\Hom_k(D_\sigma(I),k)$; see \cite[Proposition 2.5]{Jiao}. One can obtain them by writing $D_\sigma(I)$ as the cokernel of the corresponding stable Hom map, dualizing, and applying Auslander--Reiten duality to the resulting kernel.

The covariant representable resolution associated with a minimal sequence is minimal: both maps remain in the radical. Thus $\Hom_A(Y_D,-)\to D'_\sigma$ is its projective cover. The translation equivalence preserves the radical and the residue division fields, while $k$-duality interchanges socle and top. The multiplicity of $\tau I_i$ in $Y_D$ is therefore $s_i(D)$. Injective summands of $Y_D$ would split off contractible summands, and hence do not occur.
\end{proof}

Let $H=(h(I_i,I_j))_{i,j=1}^m$. We identify the split Grothendieck group with $\ZZ^m$ by the classes of the $I_i$. The matrix $H$ is invertible over $\QQ$. For a nonprojective $I_j$, let $E_j$ be the middle term of its almost split sequence. Applying $\Hom_A(I_i,-)$ shows that the column
$[\tau I_j]+[I_j]-[E_j]$ is sent by $H$ to $d_j$ times the $j$th unit vector. For projective $I_j$, use $[I_j]-[\rad I_j]$ instead: every nonisomorphism from an indecomposable to $I_j$ has image in its unique maximal submodule. Let $C$ be the matrix with these columns. Then $HC=\operatorname{diag}(d_1,\ldots,d_m)$.

For a sequence with terms $X,Y,L$ it follows that
\begin{equation}\label{eq:defdim}
 \delta_i=h(I_i,Y)+h(I_i,X)-h(I_i,L),
\end{equation}
and, for a minimal sequence of defect data $\delta,t,s$,
\begin{equation}\label{eq:middleterm}
 [L_D]=[X(t)]+[Y(s)]-H^{-1}\delta.
\end{equation}
In particular the minimal middle term is determined by those data.

\subsection{Stabilizers and the mass formula}

\begin{lemma}\label{lem:stabilizer}
If the contractible terms of $\sigma$ are $B,C$, then over a field of order $Q$,
\begin{equation}\label{eq:autseq}
 |\Aut(\sigma)|=
 Q^{h(L,Y)+h(X,L)-h(X,Y)-e(B)-e(C)}|\Aut(D_\sigma)|a_B(Q)a_C(Q).
\end{equation}
\end{lemma}
\begin{proof}
There is an exact sequence of vector spaces
\begin{align*}
0\to\Hom(X,Y)&\longrightarrow\Hom(L,Y)\oplus\Hom(X,L)\\
&\longrightarrow\End(\sigma)\longrightarrow\End(D_\sigma)\to0,
\end{align*}
where the first two maps are $u\mapsto(up,-iu)$ and
$(s,t)\mapsto(si,is+tp,pt)$. The last map is onto by lifting along the projective resolution \eqref{eq:resolution}. If a chain map induces zero on the defect, its component on $X$ lifts to $pt$ for a map $t:X\to L$. Its middle component minus $tp$ factors through $i$, giving $is$, and the chain condition gives the component $si$ on $Y$. Finally, $si=pt=is+tp=0$ implies $s=up$ and $t=-iu$. This proves exactness and
\[
 \dim\End(\sigma)=\dim\End(D_\sigma)+h(L,Y)+h(X,L)-h(X,Y).
\]
For a minimal resolution, every null-homotopic chain endomorphism has radical components. If $f$ is such an endomorphism and $g$ is arbitrary, $1-gf$ is invertible in every degree and its inverses form a chain map. Hence the kernel of $\End(\sigma_D)\to\End(D)$ lies in the Jacobson radical and the semisimple quotients agree.

The category of these finite complexes is Krull--Schmidt. Its minimal indecomposable summands correspond to the indecomposable summands of $D$. Neither type of nonzero contractible summand is isomorphic to a minimal summand or to the other contractible type. Therefore the semisimple quotient of $\End(\sigma)$ is the product of those of $\End(D)$, $\End(B)$, and $\End(C)$. Applying
\[
 |R^\times|=Q^{\dim R-\dim(R/\rad R)}|(R/\rad R)^\times|
\]
to these four endomorphism algebras proves \eqref{eq:autseq}.
\end{proof}

Define the mass, over the current extension field, by
\[
 \mathfrak M_{\delta,t,s}(E)=
 \sum_{\substack{[D]\in\operatorname{mod}\Lambda_E\\
 \delta(D)=\delta,\ t(D)=t,\ s(D)=s}}
 \frac1{|\Aut(D)|}.
\]
For fixed dimension this is a finite sum. No representation-finiteness assumption on $\Lambda$ is needed. The mass of the zero module is one. Incompatible data have mass zero.

\begin{proposition}\label{prop:mass}
Fix $X,Y,L$ and let $\delta$ be given by \eqref{eq:defdim}. If $B_s=Y\ominus Y(s)$ and $C_t=X\ominus X(t)$ denote existing direct-sum complements, then
\begin{equation}\label{eq:mass}
 F^L_{X,Y}(E)=a_L(Q)Q^{h(X,Y)-h(L,Y)-h(X,L)}
 \sum_{t,s}\frac{Q^{e(B_s)+e(C_t)}}{a_{B_s}(Q)a_{C_t}(Q)}
 \mathfrak M_{\delta,t,s}(E).
\end{equation}
The sum ranges over precisely the pairs for which both complements exist.
\end{proposition}
\begin{proof}
The group $\Aut(L_E)$ acts on the submodules counted by the left-hand side. The stabilizer of a submodule is the automorphism group of its short exact sequence, since the maps induced on the submodule and quotient are uniquely determined. For a fixed defect $D$, decomposition \eqref{eq:contracts} shows that every nonempty layer is one orbit. Its complements are determined by the endpoints. Equation~\eqref{eq:middleterm}, together with \eqref{eq:defdim}, then gives the required middle term. Apply \eqref{eq:autseq} and sum the orbit sizes.
\end{proof}

\begin{corollary}\label{cor:criterion}
If every mass $\mathfrak M_{\delta,t,s}$ is given by an element of $\R$ on all conservative extensions, then $A$ has Hall polynomials.
\end{corollary}
\begin{proof}
Every sum in \eqref{eq:mass} is finite, and all its other factors belong to $\R$. The result takes integer values on the unbounded set of conservative field sizes. Apply Lemma~\ref{lem:integer}.
\end{proof}

The mass transformation is reversible. For fixed $\delta,t,s$, form the vector in \eqref{eq:middleterm}. If it is not a nonnegative integral vector, the mass is zero. Otherwise let $L_0$ be the corresponding module and apply \eqref{eq:mass} to $X(t),Y(s),L_0$. After multiplication by $Q^{h(L_0,Y(s))+h(X(t),L_0)-h(X(t),Y(s))}/a_{L_0}(Q)$, the coefficient of $\mathfrak M_{\delta,t,s}$ is one, and all other terms have smaller $\sum_i(t_i+s_i)$. This also proves the converse of Corollary~\ref{cor:criterion}.

\section{Nilpotent cycles and square boundaries}\label{sec:cycles}
\subsection{Semilinear cyclic Hall polynomials}
Let $K$ be a finite field of order $z$, let $p\ge1$, and index vertices by $\ZZ/p\ZZ$. Fix automorphisms $\theta_i$ of $K$. A semilinear cyclic representation consists of finite-dimensional $K$-spaces $V_i$ and $\theta_i$-semilinear maps $T_i:V_i\to V_{i+1}$. A morphism is a collection of $K$-linear maps commuting with the $T_i$. We restrict to nilpotent representations: every sufficiently long consecutive composite of arrows is zero.

For $i\in\ZZ/p\ZZ$ and $l\ge1$, let $I(i,l)$ be the chain with basis positions of colours $i,i+1,\ldots,i+l-1$, with the prescribed semilinear transport between consecutive positions and zero after the last position. A label is a finitely supported family $(m_{i,l})$ of nonnegative integers, denoting $\bigoplus_{i,l}I(i,l)^{m_{i,l}}$.

\begin{lemma}\label{lem:cyclic}
For fixed chain labels, the automorphism orders and Hall numbers are polynomials in $z$ with integer coefficients, independent of the automorphisms $\theta_i$. The automorphism polynomials are monic and their inverses belong to $\ZZ[z,z^{-1},(z^j-1)^{-1}:j\ge1]$.
\end{lemma}
\begin{proof}
Put $V=\bigoplus_iV_i$ and write $T$ for the additive map on $V$ given by the arrows and set $W_h=\ker T\cap\im T^h$, degree by degree. Choose bases of the descending flags $W_h$ and lift a basis of $W_h/W_{h+1}$ through $T^h$. The resulting strings of length $h+1$ are independent: applying the highest remaining power of $T$ to a linear relation gives a relation among the chosen socle basis vectors. Semilinearity applies an invertible field automorphism to each coefficient, so the coefficients vanish. The number of chain positions is
\[
 \sum_{h\ge0}(h+1)\dim_K(W_h/W_{h+1})
 =\sum_{h\ge0}\dim_KW_h=\dim_KV.
\]
The last identity follows from
$\dim W_h=\dim\im T^h-\dim\im T^{h+1}$.
This proves existence of the chain decomposition. Its multiplicities are determined by the colours and dimensions of $W_h/W_{h+1}$, proving uniqueness.

Put $r_{i,h}(V)=\dim_K(\im T^h)_i$ and
$z_{i,h}(V)=r_{i,h}(V)-r_{i,h+1}(V)$. The chain multiplicities can equivalently be recovered by
\begin{equation}\label{eq:chainrecover}
 m_{i,l}=z_{i+l-1,l-1}(V)-z_{i+l,l}(V),
\end{equation}
with colour indices read modulo $p$. Thus fixed total dimension gives a finite label set.

For integers $0\le e\le b$, the Gaussian polynomial
$\gauss be z=\prod_{j=0}^{e-1}(z^{b-j}-1)/(z^{e-j}-1)$
counts $e$-dimensional subspaces of $K^b$ and is zero outside this range. 
First count a semisimple submodule $S$ of $L$ with quotient $X$. Write $s_i=\dim_KS_i$ and
\[
 W_{i,h}=(\soc L)_i\cap(\im T^h)_i,\qquad
 c_{i,h}=r_{i,h}(L)-r_{i,h}(X).
\]
A submodule $U\simeq S$ is a tuple of subspaces $U_i\subseteq(\soc L)_i$. Since
$\im(T^h:L/U\to L/U)=(\im T^h+U)/U$,
its quotient is $X$ precisely when
$\dim(U_i\cap W_{i,h})=c_{i,h}$ for all $i,h$, by \eqref{eq:chainrecover}. Define
\[
 b_{i,h}=\dim W_{i,h}-\dim W_{i,h+1},\qquad
 e_{i,h}=c_{i,h}-c_{i,h+1}.
\]
The count is zero unless $c_{i,0}=s_i$, the tails vanish, and $0\le e_{i,h}\le b_{i,h}$. Otherwise it is
\begin{equation}\label{eq:semisimplecount}
 F^L_{X,S}(z)=
 \prod_i\left(z^{\sum_{h<a}e_{i,h}(b_{i,a}-e_{i,a})}
 \prod_{h\ge0}\gauss{b_{i,h}}{e_{i,h}}{z}\right).
\end{equation}
To obtain \eqref{eq:semisimplecount}, choose the image in each successive flag quotient and then its graph into the deeper unchosen quotient spaces. These choices contribute, respectively, the Gaussian factors and the displayed power of $z$.

Use the unmodified Hall multiplication $u_Xu_Y=\sum_LF^L_{X,Y}u_L$. For a module $M$, choose $l\ge1$ with $\rad^lM=0$, and let $S_h=\rad^hM/\rad^{h+1}M$, where the radical is the arrow image. Expanding $u_{S_0}\cdots u_{S_{l-1}}$ by left-associated multiplication uses only \eqref{eq:semisimplecount}. A term $u_N$ in this product corresponds to a descending flag with these semisimple quotients. It satisfies
$\rad^hN\subseteq N_h$ and hence $r_{i,h}(N)\le r_{i,h}(M)$. Equality for all $i,h$ implies $N\simeq M$ by \eqref{eq:chainrecover}; then the only possible flag is the radical filtration. Its coefficient is one. Therefore
\[
 u_{S_0}\cdots u_{S_{l-1}}
 =u_M+\sum_{\sum_{i,h\ge1}r_{i,h}(N)<\sum_{i,h\ge1}r_{i,h}(M)}a_{M,N}(z)u_N
\]
with integer polynomial coefficients. Induction expresses every $u_M$ as an integer polynomial combination of words in semisimple modules. Expand the second factor of $u_Xu_Y$ in such words and multiply successively on the right by semisimple factors. This proves the Hall assertion, with the same polynomials for every holonomy.

A map between two chains is determined by the image of its first generator: this image has the required colour and is killed by the prescribed power of $T$. Its free coordinates are counted by a power of $z$. The endomorphism ring of one chain has residue field $K$; the leading generator coefficient is arbitrary in $K$, and zero leading coefficient gives the nilpotent radical. Thus, if $M=\bigoplus_I I^{m_I}$ and $|\End(M)|=z^{e_M}$,
\[
 |\Aut(M)|=z^{e_M-\sum_I m_I^2}\prod_I|\operatorname{GL}_{m_I}(K)|.
\]
The exponent and factors depend only on the labels, proving the final assertion.
\end{proof}

A Hall flag polynomial with successive quotients $T_1,\ldots,T_n$ means the coefficient of $u_L$ in $u_{T_n}\cdots u_{T_1}$, or equivalently the number of flags $0=L_0\subseteq\cdots\subseteq L_n=L$ with $L_j/L_{j-1}\simeq T_j$. Lemma~\ref{lem:cyclic} applies as well to a disjoint union of semilinear cycles, by taking products.

\subsection{Total mass at a square boundary}\label{sec:square}
For integers $n,s\ge1$, consider vertices $(i,t)$ with $0\le i<n$ and $t\in\ZZ/s\ZZ$. Put a copy of $K$ at each vertex. The forward arrows
$a_{i,t}:(i,t)\to(i+1,t)$ are linear, the central arrows
$\varepsilon_t:(0,t)\to(0,t+1)$ are $\theta_t$-semilinear, and the reverse arrows
$b_{i,t}:(i+1,t)\to(i,t+2)$ have transport $\theta_{t+1}\theta_t$.
Composition is written with the rightmost map acting first. For $n\ge2$ impose
\begin{align}
 \varepsilon_{t+1}\varepsilon_t&=b_{0,t}a_{0,t},\label{eq:square}\\
 a_{i-1,t+2}b_{i-1,t}&=b_{i,t}a_{i,t}\quad(1\le i<n-1),\notag\\
 a_{n-2,t+2}b_{n-2,t}&=0.\notag
\end{align}
For $n=1$ impose just $\varepsilon_{t+1}\varepsilon_t=0$.

\begin{proposition}\label{prop:squaretotal}
For every fixed vertex-dimension vector, the total mass $\sum_{[V]}1/|\Aut(V)|$ of these representations is given by an element of $\ZZ[z,z^{-1},(z^j-1)^{-1}:j\ge1]$, independent of the $\theta_t$.
\end{proposition}
\begin{proof}
Fix the forward-chain representations $M_t$ of $0\to1\to\cdots\to n-1$. They are sums of interval representations, so only finitely many types occur in the fixed dimension. At the initial vertex define
\[
 F_{0,t}=0,\quad F_{j,t}=\ker(M_{0,t}\to M_{j,t})\ (1\le j<n),\quad F_{n,t}=M_{0,t}.
\]
Write $f_{j,t}=\dim_KF_{j,t}$. The image of evaluation at zero on $\Hom(M_{t+2},M_t)$ consists exactly of maps preserving these endpoint flags, with the indicated scalar transport. Indeed, an interval starting at zero and ending at $a$ contributes a flag direction born in layer $a+1$; a map between two such directions exists precisely when the target birth layer is no larger than the source birth layer. Intervals starting later contribute zero at the endpoint. This proves both inclusion and surjectivity onto every allowed block.

The transpose, under trace pairings, of the reverse-arrow relation map is the usual chain Hom differential. Its cokernel is $\Ext^1(M_{t+2},M_t)$, with the source transported to make the maps linear. Pairing the boundary term in \eqref{eq:square} with the kernel of this differential shows that reverse arrows exist precisely when
\begin{equation}\label{eq:lowerflags}
 \varepsilon_{t+1}\varepsilon_t F_{j,t}\subseteq F_{j-1,t+2}
 \quad\text{for all }t,j.
\end{equation}
In adapted flag bases the Hom evaluation has a block from source layer $a$ to target layer $b$ exactly when $b\le a$. Its trace annihilator has only the opposite, strictly decreasing blocks, which is \eqref{eq:lowerflags}. If this condition holds, the reverse-arrow solution set is an affine space of size
$z^{\sum_t\dim_K\Ext^1(M_{t+2},M_t)}$.
Transport does not change an interval label, so these dimensions depend only on the forward-chain types.

Condition~\eqref{eq:lowerflags} forces every composite of $2n$ central arrows to vanish. Let $N$ range over the finite set of nilpotent cyclic chain labels with vertex dimensions $f_{n,t}$. Define $\Phi N$ by replacing each central arrow by the two-step map from $t$ to $t+2$. It is a representation of $\gcd(s,2)$ disjoint cycles. On every chain, taking alternate positions gives the chain decomposition of $\Phi N$; this operation on labels is independent of the field and of its automorphisms.

Let $H_{\Phi N,f}(z)$ count flags of subrepresentations of $\Phi N$ whose $j$th layer has vertex dimensions $f_{j,t}-f_{j-1,t}$ and is semisimple. Equivalently these flags satisfy \eqref{eq:lowerflags}. It is a Hall flag polynomial by Lemma~\ref{lem:cyclic}. The order of the stabilizer of the given flags, in $\prod_t\operatorname{GL}(M_{0,t})$, is
\[
 p_f(z)=\prod_t\left(
 z^{\sum_{a<b}(f_{a,t}-f_{a-1,t})(f_{b,t}-f_{b-1,t})}
 \prod_{j=1}^n|\operatorname{GL}_{f_{j,t}-f_{j-1,t}}(K)|\right).
\]
Double-count pairs consisting of a central representation of type $N$ and an admissible flag. Its representation orbit has size $\prod_t|\operatorname{GL}(M_{0,t})|/a_N(z)$; the space of flags has size $\prod_t|\operatorname{GL}(M_{0,t})|/p_f(z)$. Hence the number of central-arrow choices for the fixed flags is
\[
 p_f(z)\sum_N\frac{H_{\Phi N,f}(z)}{a_N(z)}.
\]
Finally quotient by the automorphisms of the fixed forward chains. The total mass is
\begin{equation}\label{eq:coretotal}
 \sum_{[M_t]}
 \frac{z^{\sum_t\dim_K\Ext^1(M_{t+2},M_t)}}{\prod_t a_{M_t}(z)}
 \ p_{f(M)}(z)\sum_N\frac{H_{\Phi N,f(M)}(z)}{a_N(z)}.
\end{equation}
All sums are finite. Interval automorphism orders and the $a_N$ have only powers of $z$ and $z^j-1$ as factors. This proves the result, including $n=1$, for which the single flag step requires the two-step maps to vanish.
\end{proof}

\section{Masses on modulated meshes}\label{sec:counting}
\subsection{Linear relations and Gaussian moments}
In this section, $h(U,V)$ denotes the dimension over the base field of the Hom space in the category containing $U,V$; the category will always be specified.

A finite species consists of vertex fields $K_i/k$ and finite-dimensional $(K_j,K_i)$-bimodules $B_a$ for arrows $a:i\to j$. Its tensor algebra is $H=T_{\prod_iK_i}(\bigoplus_aB_a)$. We use covariant representations: an arrow acts by a $K_j$-linear map $B_a\otimes_{K_i}V_i\to V_j$. This convention will be applied to the opposite of the stable mesh category when counting right defect modules.

Let $F$ be a species automorphism, including a vertex permutation $\rho$, field isomorphisms along that permutation, and invertible compatible transport of the arrow bimodules. It induces an exact equivalence, also denoted $F$, on representations. In the split case our convention is $(FV)_i=V_{\rho i}$, with transported arrow maps. For a dimension vector $v$, write $Fv$ for its transported dimension vector. For dimension vectors measured over the vertex fields define
\begin{equation}\label{eq:euler}
 \chi_H(v,w)=\sum_i[K_i:k]v_iw_i
 -\sum_{a:i\to j}(\dim_kB_a)v_iw_j.
\end{equation}
For representations $U,V$, the two-term complex
\begin{equation}\label{eq:homcomplex}
 \begin{gathered}
 \bigoplus_i\Hom_{K_i}(U_i,V_i)\xrightarrow{d}
 \bigoplus_{a:i\to j}\Hom_{K_j}(B_a\otimes_{K_i}U_i,V_j),\\
 d(f)_a=f_jU_a-V_a(1\otimes f_i)
 \end{gathered}
\end{equation}
has kernel $\Hom_H(U,V)$ and cokernel $\Ext_H^1(U,V)$. For the latter assertion, write an extension on the direct-sum vertex spaces: its off-diagonal arrow blocks are arbitrary, and changing a vertex splitting adds $d(f)$. The difference of the dimensions of the two spaces is \eqref{eq:euler}.

The \emph{transpose mesh} adds trace-dual reverse arrows and imposes the transpose of \eqref{eq:homcomplex}. The dual of a bimodule is $D_kB_a=\Hom_k(B_a,k)$, with reversed actions and the target field transported by $F$. The field trace identifies this with either relative dual; the trace pairing is nondegenerate because finite fields are separable. These choices specify the relations independently of bases. In split coordinates an old arrow $a:i\to j$ has a reverse arrow $y_a:\rho j\to i$, and the relation at $i$ is
\[
 \mu_i(x,y)=\sum_{t(a)=i}x_a y_a-\sum_{s(a)=i}y_a x_{\rho a}=0.
\]
Here $s(a)$ and $t(a)$ denote source and target.

We also allow disjoint special cycles of vertices $v_t$, indexed by $t\in\ZZ/s\ZZ$ for a cycle of length $s\ge1$, with $\rho v_t=v_{t-2}$ and arrows $\varepsilon_t:v_t\to v_{t+1}$. At $v_{t+2}$ the relation becomes
$\mu_{v_{t+2}}+\varepsilon_{t+1}\varepsilon_t=0$.
The central-arrow bimodules are invertible and their twofold transport equals that of the relation. Scalar signs may be absorbed in the reverse arrows.

Fix a dimension vector $v$ and nonnegative vectors $r,u$. For a representation $D$, the space $(\soc_{\rm ar}D)_i$ is the common kernel of arrows leaving $i$, and $(\rad_{\rm ar}D)_i$ is the sum of images of arrows entering $i$. Set $\topm_{\rm ar}D=D/\rad_{\rm ar}D$. These spaces are taken vertexwise; all dimensions in the following formula are over the vertex fields. Put
\begin{equation}\label{eq:moment}
 Z_{v;r,u}(E)=\sum_{[D],\,\dim D=v}\frac1{|\Aut D|}
 \prod_i\gauss{\dim(\soc_{\rm ar}D)_i}{r_i}{Q^{d_i}}
 \gauss{\dim(\topm_{\rm ar}D)_i}{u_i}{Q^{d_i}},
 \qquad Q=|E|.
\end{equation}
Here $d_i=[K_i:k]$, $E$ preserves all vertex fields, and vertex dimensions are over $K_i\otimes_kE$. This is the mass of triples $(D,S,K)$ with $S\subseteq\soc_{\rm ar}D$, $K\supseteq\rad_{\rm ar}D$, $\dim S=r$, and $\dim(D/K)=u$. The two subspaces are chosen independently.

Fix the old-arrow representation $M$ and its subspaces $S,K$, and put $X=M/S$. Write $\alpha:K\hookrightarrow M$, $\beta:M\twoheadrightarrow X$, and $h=\beta\alpha$. A reverse arrow killing $S$ and landing in $K$ is uniquely a compressed arrow from the corresponding transported $X$-space to a $K$-space. Its parameter space is the dual of the degree-one space of \eqref{eq:homcomplex} for $(K,FX)$; the relation space is the dual of the degree-zero space. Denote these dual spaces by $C^1,C^0$. The relation map $\mu:C^1\to C^0$ has transpose $d$.

On a special cycle write
\[
 \varepsilon_t=\alpha_{t+1}z_t\beta_t,\qquad
 z_t:X_t\to K_{t+1},\qquad c_t(z)=z_{t+1}h_{t+1}z_t,
\]
with the prescribed transports. At ordinary vertices put $c(z)=0$. The compressed equations are $\mu(y)+c(z)=0$. For a nontrivial additive character $\psi:E\to\mathbb C^\times$, orthogonality gives
\begin{equation}\label{eq:fourierexact}
 \#\{(y,z):\mu(y)+c(z)=0\}
 =Q^{\dim_EC^1-\dim_EC^0}
 \sum_{f\in\ker d}\sum_z\psi\langle f,c(z)\rangle .
\end{equation}
Indeed, insert $Q^{-\dim C^0}\sum_{f\in(C^0)^*}\psi\langle f,\mu(y)+c(z)\rangle$ and sum over $y$. The sum vanishes unless $d(f)=0$, when it contributes $Q^{\dim C^1}$. The pairing includes field traces, and
$\dim C^1-\dim C^0=-\chi_H(v-u,F(v-r))$.

\subsection{String and band representations}
On each orbit $o$ of $\rho$, the vertexwise data $(K,M,X,f)$ form a diagram
\begin{equation}\label{eq:local}
 K_t\xrightarrow{\alpha_t}M_t\xrightarrow{\beta_t}X_t,
 \qquad K_t\xrightarrow{f_t}X_{t-1},
\end{equation}
after successively identifying the vertex fields. The maps $\alpha_t$ are injective and $\beta_t$ are surjective. An orbit of length $p$ gives a semilinear cycle with $3p$ vertices, with orientation $K\to M\to X\leftarrow K$ and residual field transport on one edge.

The string and band classification applies to this local category. Its underlying quiver is a cycle, has no special loops or relations, has at most two incoming and outgoing arrows at each vertex, and each arrow has at most one composable successor and predecessor. Thus the Main Theorem of \cite{BTCB} applies; its hypotheses on special-loop polynomials are vacuous. Strings are finite intervals in the universal covering chain, pushed down with the prescribed transports. A closed reduced walk traverses the entire cycle, so every arrow of a band is invertible. Bands correspond to finite-dimensional modules over a skew Laurent ring. Every local object therefore decomposes uniquely as
\[
 E_o=W_o\oplus B_o,
\]
where $W_o$ is a sum of finite strings and $B_o$ has all arrows invertible. Write $b_o$ for its constant vertex rank, and regard $b=(b_o)_o$ as a vertex-dimension vector by assigning $b_o$ to every vertex in the orbit $o$. Injectivity of $\alpha$ and surjectivity of $\beta$ are tested on each summand; every band satisfies them. Fixed dimensions allow only finitely many string labels and multiplicities.

For a string, a diagonal endomorphism coefficient is arbitrary in its vertex field and is transported along the finite interval. Its residue field is that vertex field. Hom spaces between strings are found by propagating coefficients along overlap intervals and imposing zero at obstructed ends. Hence their dimensions and automorphism polynomials $a_{W_o}(Q)$ are unchanged by vertex-conservative extension, and $a_{W_o}^{-1}\in\R$.

We describe all old-arrow choices for fixed local diagrams. In the split case, an old-arrow orbit of length $l$ has endpoint periods dividing $l$. Pull both endpoint diagrams to the common $3l$-vertex cycle. The arrow orbit is a morphism between these pullbacks whose $M$-component kills the source $\ker\beta$ and lands in the target $\im\alpha$. Denote this vector space by $H_a^0(E_o,E_{o'})$ and its dimension over the current base field by $h_a^0$. Its induced $K$- and $X$-components are unique, and the $f$-square is exactly $d(f)=0$. With modulation, tensor the source with the arrow bimodule and use its invertible transport in the same equations; this defines $H_a^0$ without choosing an invariant arrow basis.

For local objects $E,E'$, let $\Hom^0(E,E')$ be the space of morphisms whose $M$-components kill $\ker\beta$ and land in $\im\alpha'$ of the target. Let $S(E)$ have value $\ker\beta$ at its $M$-vertices and zero elsewhere, and let $E^+$ be the kernel of the quotient supported on the $M$-vertices with value $\coker\alpha$. Then
\[
 \Hom^0(E,E')=\Hom(E/S(E),(E')^+).
\]
These two operations replace strings by sums of intervals and leave bands unchanged.

\begin{lemma}\label{lem:mixedhom}
The string--string, string--band and band--string dimensions in every $H_a^0$, and the corresponding local Hom dimensions, depend only on string labels and band ranks.
\end{lemma}
\begin{proof}
A pullback of a finite interval through a cyclic cover is a union of finite intervals: a starting lift determines the lift of its entire coefficient interval. For a map from an interval to a band, attach a target vector to each coefficient position. Interior equations transport it invertibly; an endpoint either leaves it free or sets it to zero. Each surviving interval contributes one full fibre. Maps in the opposite direction use row vectors. For two intervals, use the connected overlap intervals.

With an arrow bimodule, each coefficient pair carries the corresponding vector fibre. Its holonomy gives invertible maps between successive fibres. Since their coefficient interval has distinct endpoints, these maps only transport a free fibre. The dimension is the sum of the dimensions of surviving fibres. The operations $E/S(E)$ and $(E')^+$ delete coefficient positions, so the same calculation applies to the shortened intervals. This proves the assertion for different endpoint periods and for semilinear transport.
\end{proof}

The mass for fixed local diagrams and all old-arrow choices is
$Q^{\sum_a h_a^0(E_o,E_{o'})}/\prod_o|\Aut E_o|$.
To see this, fix the vertex spaces and divide by their change-of-basis groups. Adding frames for $K,X$ introduces exactly the groups removed by the injective and surjective maps in \eqref{eq:local}. Each old-arrow orbit is then the independent vector space just described.

\subsection{The quadratic character sum}
On a special cycle the phase in \eqref{eq:fourierexact} is
\begin{equation}\label{eq:phase}
 \sum_t\Tr(f_tz_{t+1}h_{t+1}z_t),\qquad f_t:K_{t+2}\to X_t,
\end{equation}
where the matrix trace is followed by the field trace to $E$. For a fixed local diagram $E_{\rm loc}$, let $\mathcal W(E_{\rm loc})$ be the sum over all compressed central maps $z$ of $\psi$ applied to this phase. If there are no special arrows, set $\mathcal W(E_{\rm loc})=1$.

\begin{lemma}\label{lem:pathsum}
For fixed local strings $W$ and band ranks $b$, summation of $\psi$ applied to \eqref{eq:phase} factors as
\[
 \mathcal W(W\oplus B)=Q^{\gamma(W,b)}\mathcal W(B),
\]
where $\gamma(W,b)$ is a nonnegative integer independent of band parameters and of $\psi$.
\end{lemma}
\begin{proof}
Use coefficient bases on strings. Their $h$ and $f$ maps are partial matchings with the prescribed field transport. Index a coefficient of $z_t$ by its target position $k$ in $K_{t+1}$ and source position $x$ in $X_t$. Its successor in \eqref{eq:phase}, when defined, is
\[
 (k,x)\longmapsto(f_t^{-1}x,h_{t+1}k).
\]
Here the inverse is that of a partial matching. Each coefficient has at most one predecessor and successor. Two successor steps move both coefficient positions forward on their local $h,f$ covering chains. A position on a finite string reaches an endpoint, so every component containing a string is a finite path.

For a path with $m$ scalar variables and nonzero edge coefficients,
\[
 \sum_{x_1,\ldots,x_m\in E}
 \psi\left(\sum_{j=1}^{m-1}c_jx_jx_{j+1}\right)=Q^{\lceil m/2\rceil}.
\]
Summing the first variable sets the second to zero and contributes $Q$; repeat at the new endpoint. An isolated variable also contributes $Q$. When one position belongs to a band, retain its entire vector space. The same path carries constant-rank vector fibres with nondegenerate pairings on adjacent fibres, and the same orthogonality calculation gives the appropriate power of the field size. Semilinear transport is bijective and field trace is nondegenerate, so this remains valid over the original field. The sum of the fibre dimensions times $\lceil m/2\rceil$ is $\gamma(W,b)$. The unintegrated variables form exactly the band--band block.
\end{proof}

\subsection{Band masses and Gaussian inversion}
Let $J$ be the vertex support allowed for the representations and let $P$ be the union of the $\rho$-cycles entirely contained in $J$. Assume that special cycles are contained in $P$. Every nonzero band has equal positive rank along its orbit, so bands are supported on $P$.

String and band summands have disjoint indecomposable types. Taking semisimple quotients of their endomorphism rings gives
\[
 |\Aut(W_o\oplus B_o)|=a_{W_o}(Q)|\Aut B_o|
 Q^{h(W_o,B_o)+h(B_o,W_o)}.
\]
For fixed $W,b$ define
\begin{align*}
 \eta(W,b)={}&\sum_{a:o\to o'}\bigl(
 h_a^0(W_o,W_{o'})+h_a^0(W_o,B_{o'})+h_a^0(B_o,W_{o'})\bigr)\\
 &-\sum_o\bigl(h(W_o,B_o)+h(B_o,W_o)\bigr).
\end{align*}
Here all dimensions are over the current base field. Lemma~\ref{lem:mixedhom} makes this independent of the chosen band types. The remaining band sum is
\[
 P_b(Q;\psi)=\sum_{[B],\,\operatorname{rank}B=b}
 \frac{Q^{\sum_{a\text{ in }P}h_a^0(B_o,B_{o'})}}{\prod_o|\Aut B_o|}
 \mathcal W(B).
\]
Without special arrows the character factor is one. Combining the preceding computations yields
\begin{equation}\label{eq:mixed}
 Z_{v;r,u}(Q)=Q^{-\chi_H(v-u,F(v-r))}
 \sum_b\sum_W\frac{Q^{\eta(W,b)+\gamma(W,b)}}{\prod_o a_{W_o}(Q)}
 P_b(Q;\psi).
\end{equation}
The $K,M,X$ dimensions of $W$ must be $v-u-b,v-b,v-r-b$, respectively; these conditions specify the finite sums. At ordinary orbits $\gamma=0$.

\begin{proposition}\label{prop:recovery}
Suppose that the total representation mass on the full subquiver $P$, for each $\rho$-invariant dimension vector, is given by an element of $\R$ on all vertex-conservative extensions. Then every fixed-dimension, fixed-arrow-top and fixed-arrow-socle mass on $J$ has this property.
\end{proposition}
\begin{proof}
Apply \eqref{eq:mixed} on $P$ with $r=u=0$ and $\rho$-invariant $v$. Denote its finite sum of string coefficients at rank $b$ by $A_{v,b}(Q)\in\R$. At $b=v$ the only string object is zero, with coefficient one. Hence
\begin{equation}\label{eq:bandrecursion}
 P_v(Q;\psi)=Q^{\chi_H(v,Fv)}Z^P_{v;0,0}(Q)
 -\sum_{b<v}A_{v,b}(Q)P_b(Q;\psi).
\end{equation}
Here $b<v$ means componentwise inequality with at least one strict component. Starting with $P_0=1$, induction on total dimension gives $P_v\in\R$; the recursion also gives independence of $\psi$. Multiple periodic components can be treated as a direct product.

Substitute into \eqref{eq:mixed} to obtain every Gaussian moment in $\R$. For fixed $v$, top and socle dimensions are bounded. Each matrix
$\bigl(\gauss sr{Q^{d_i}}\bigr)_{0\le r,s\le v_i}$
is triangular with diagonal entries one, and its inverse has entries in $\ZZ[Q]$ by finite back-substitution. Invert in both sets of variables in \eqref{eq:moment}.
\end{proof}

\subsection{Representation masses for Dynkin species}
For an ordinary transpose mesh with representation-finite species $H$, fix an old-arrow representation $M$. The reverse-arrow space is $\Ext_H^1(M,FM)^*$ by \eqref{eq:homcomplex}. Its total mass is therefore
\begin{equation}\label{eq:ordinarytotal}
 \sum_{[M],\,\dim M=v}\frac{Q^{\dim_E\Ext_H^1(M,FM)}}{a_M(Q)}.
\end{equation}
On conservative extensions of $H$, its labels, Hom dimensions and the permutation of labels induced by $F$ are fixed. Thus \eqref{eq:ordinarytotal} is a finite sum in $\R$.

\begin{lemma}\label{lem:residues}
If $H$ is acyclic and its symmetrized Euler form is positive definite, every indecomposable $H$-module has endomorphism algebra isomorphic as a $k$-algebra to a vertex field. Its conservative period is the least common multiple of the vertex-field degrees.
\end{lemma}
\begin{proof}
At a sink $i$, a representation without the simple summand at $i$ has a surjective map
$\bigoplus_{a:j\to i}B_a\otimes_{K_j}M_j\to M_i$.
Replace $M_i$ by its kernel and reverse the incident bimodules using relative duals. The inverse at the new source takes the cokernel. The two exact sequences give mutually inverse functors on the subcategories without the corresponding simple summands. They preserve indecomposability and endomorphism rings.

Let $E_H$ be the Euler matrix in a topological order and $G_H=E_H+E_H^{\mathsf t}$. A sink reflection changes coordinate $i$ to
$-m_i+\sum_{a:j\to i}(\dim_kB_a)m_j/[K_i:k]$.
Solving these coordinate updates successively along a complete sink sequence gives $c=-E_H^{-1}E_H^{\mathsf t}$. The transformations preserve the integral lattice and $G_H$. Positive definiteness implies that the integral isometry group is finite: each column has fixed bounded $G_H$-length and hence finitely many possible integral values. Thus $c^h=1$ for some $h$. Also $c-1=-E_H^{-1}G_H$ is invertible, so $1+c+\cdots+c^{h-1}=0$.

Repeat the sink sequence on a nonzero indecomposable until it becomes the simple at the current sink. If this never happens, $m,cm,\ldots,c^{h-1}m$ are all nonnegative, the first is nonzero, and their sum is zero, a contradiction. Endomorphism rings are preserved along these reflections and the final simple has a vertex field as its endomorphism ring. The same bound gives finitely many possible indecomposables. Every vertex field already occurs at its simple, giving the asserted least common multiple.
\end{proof}

\section{Standard forms and finite completion}\label{sec:structure}
\subsection{Descent of the mesh presentation}

We first explain why geometric standardness gives a presentation over the original field. For a finite-dimensional algebra $C$, its radical grading is
\[
 \gr C=\bigoplus_{j\geq0}\rad^j C/\rad^{j+1}C.
\]
A radical-graded algebra is an algebra isomorphic to this graded algebra, with its grading forgotten.

\begin{lemma}\label{lem:grading}
Let $k$ be a finite field. If $C_{\bar k}$ is radical-graded, then $C$ is radical-graded.
\end{lemma}
\begin{proof}
Put $B=\gr C$. Scalar extension commutes with the radical because $k$ is perfect. Choose a geometric isomorphism $C_{\bar k}\to B_{\bar k}$. Its associated graded map is a graded automorphism of $B_{\bar k}$; composing with its inverse makes the associated graded map the identity.

Let $U$ be the reduced algebraic group of automorphisms of $B$ inducing the identity on its associated graded algebra. This group is smooth since $k$ is perfect. It is connected: write an automorphism in homogeneous blocks. Its diagonal blocks are identities and its other blocks go from degree $a$ to a strictly larger degree $b$. Conjugation by the grading action multiplies these blocks by $t^{b-a}$. It extends to $t=0$, with value the identity. Thus every geometric point of $U$ belongs to the identity component. The isomorphism scheme of maps $C\to B$ inducing the graded identity is geometrically nonempty. Choosing one geometric isomorphism identifies it, by composition, with the automorphism scheme inducing the graded identity. Reduction commutes with extension of the perfect field $k$, so its reduced scheme becomes $U_{\bar k}$ with its translation action. It is therefore a smooth $U$-torsor over $k$. Lang's theorem \cite{Lang} gives a rational point of this torsor.
\end{proof}


Take one representative of each indecomposable $A$-module and let $G$ be their direct sum. Over $\bar k$, $G_{\bar k}$ contains every geometric indecomposable. Indeed, any geometric indecomposable is defined over a finite separable extension $E/k$, and is a summand of the scalar extension of its restriction to $k$. Also
\[
 \End_A(G)/\rad\End_A(G)=\prod_i\FF_{q^{d_i}},
\]
so after splitting the fields its primitive summands have multiplicity one. Thus $\End_A(G)_{\bar k}$ is the basic endomorphism algebra of all geometric indecomposables.

In a mesh category, arrows are irreducible morphisms and the defining relations are homogeneous of degree two. In finite representation type its arrow ideal is nilpotent, so the grading is the radical grading. Lemma~\ref{lem:grading} descends this grading to $\End_A(G)$. One can align the vertex idempotents with the original summand projections: if two orthogonal systems $(e_i)$ and $(f_i)$ coincide modulo the radical, then
$u=\sum_i f_i e_i$ lies in $1+\rad\End_A(G)$ and satisfies $f_i u=u e_i$.

Here is the resulting modulated presentation. A partial translation quiver has a partially defined injective translation $\tau$ on vertices. For every arrow $a:y\to x$ with $\tau x$ defined there is a companion
\[
 \sigma a:\tau x\longrightarrow y;
\]
There is also an inverse companion $\sigma^{-1}a:x\to\tau^{-1}y$ whenever $\tau^{-1}y$ is defined. These partial operations are mutually inverse, and $\sigma^2$ is translation on arrows wherever defined. The companion pairings are nondegenerate. These assertions hold in the Auslander--Reiten quiver and remain true after deleting projectives whenever all the indicated endpoints remain. Over a splitting field, the relation at $x$ is the sum of the length-two paths $a\,\sigma a$, with the coefficients from its almost-split sequence. Over $k$, the vertices have their residue fields, the arrow spaces are the corresponding irreducible-morphism bimodules, and the relations are the coevaluation tensors for the almost-split pairings. The geometric quadratic presentation implies the presentation over $k$: scalar extension is faithfully flat, so the descended degree-two relations generate the entire kernel.

The stable Auslander algebra
\[
 \Lambda=\End_A(G)/[\operatorname{proj} A]
\]
is obtained by deleting the projective vertices and every path through them. A map between nonprojective indecomposables factoring through a projective lies in the square of the radical, since neither factor can split. Thus this quotient does not change the remaining arrow spaces. Its relation at an old vertex is exactly the sum of the remaining length-two mesh paths. The sum can be empty. Translation is retained only when both translated endpoints remain.

\subsection{Periodic components}

Call a vertex periodic if every positive and negative iterate of its translation exists. The full subquiver on periodic vertices of a representation-finite Auslander--Reiten quiver is a disjoint union of quotients of repetitive Dynkin quivers; this is the stable-subquiver statement of Riedtmann's structure theorem \cite[Hauptsatz]{Riedtmann}. It applies to the stable full subquiver of an arbitrary representation-finite algebra. Its vertices need not constitute all nonprojective vertices.

We record the cuts and their finite-field forms needed for counting. For Dynkin trees other than $A_{2n}$, orient the tree away from its central fixed vertex, or choose an orientation invariant under its diagram automorphisms. The automorphisms of the repetitive quiver are generated by translation and diagram automorphisms preserving such a cut. For $A_{2n}$ there is instead a generator $\rho$ with $\rho^2=\tau$. These assertions, including the exceptional generator, are given explicitly in \cite[Proposition 10]{Bongartz23}. 

In the first case, the forward arrows form a finite union of Dynkin species $H$; the companions are the transpose arrows. Deck transformations and finite-field descent preserve this division. The automorphism transporting one mesh to the next becomes the species automorphism $F$ of Section~\ref{sec:counting}. Folding a positive definite Dynkin form remains positive definite: the folded form is its restriction to the invariant subspace, with the vertex-field degrees as symmetrizing weights. Lemma~\ref{lem:residues} and \eqref{eq:ordinarytotal} therefore give the total mass on every vertex-conservative extension.

The exceptional cut has the following coordinates. Label the two arms of $A_{2n}$ by $(i,p)$, where $0\leq i<n$ is distance from the centre and $p=0,1$ distinguishes the arms. If $l\in\ZZ$ is the repetition coordinate, set $t=2l+p$. Then $\tau$ changes $t$ by two and $\rho$ changes it by one, after choosing the direction of indexing. The outward arm arrows are $a_{i,t}$; the central arrows are $\varepsilon_t$. Their companions are the reverse arm arrows $b_{i,t}$. Replacing the reverse arrows at distance $i$ by $(-1)^{i+1}$ times themselves gives precisely the square-boundary equations of Section~\ref{sec:square}. A quotient by a power of $\rho$ makes $t$ cyclic. The equality of the residue-field degrees follows from the action on these coordinates. Suppose Frobenius permutes $r$ geometric components transitively, and write one as $\mathbb Z A_{2n}/\langle\rho^s\rangle$. Its $r$th power lifts to a power $\rho^a$, by the automorphism description above. At every distance $i$ from the centre it acts on $t\in\mathbb Z/s\mathbb Z$ by $t\mapsto t+a$. Every geometric vertex therefore has Frobenius orbit length $rs/\gcd(s,a)$. This length is the degree of its descended residue field, independently of $i$.

The same description holds without splitting the field. Along an arm, choose successive identifications of the vertex fields and trivializations of the invertible arrow bimodules. There is no obstruction, since an arm is a tree. At the centre, retain the invertible bimodule of each $\varepsilon_t$ and its field transport. The two summands in the central relation are the tensor product of two successive central bimodules and the first arm-return bimodule. Each projection of the relation bimodule is an isomorphism: it is one after extension to $\bar k$, hence one before extension. The internal relations identify consecutive arm-return bimodules in the same way. Choose the reverse-arrow identifications successively along the arm. This gives the displayed square-boundary equations over the original fields, with all remaining transport in the $\varepsilon_t$. Proposition~\ref{prop:squaretotal} applies to the resulting semilinear central cycle.

\subsection{Completion of the finite translation intervals}

We complete the finite partial quiver underlying $\Lambda$ by adding vertices.

\begin{lemma}\label{lem:completion}
A finite modulated partial mesh quiver of the preceding form is a full vertex quotient of a finite transpose mesh with square-boundary components. Its original periodic subquiver is unchanged. Every added translation cycle has a new vertex.
\end{lemma}
\begin{proof}
Every component of the partial permutation $\tau$ is a cycle or an interval. Give an interval coordinates $0,\ldots,l_i$, with $\tau$ decreasing the coordinate by one. Let $L=\max_i l_i$, taking $L=0$ if there are no intervals. For each old translation cycle $c$, let $m_c$ be its length and let $\alpha_c$ be the automorphism of its starting vertex field obtained by composing the prescribed field transports once around the cycle. Write $h_c$ for the finite order of $\alpha_c$. Choose an even integer $M>2L+2$ divisible by every $m_ch_c$. Then transport through $M$ steps is the identity on each old cycle. Complete each interval to a cycle of length $M$.

Follow companion arrows. A chain that is already closed has both endpoints on old periodic cycles and is left unchanged. Otherwise, if one arrow is
$a_0:(i,p)\to(j,q)$, its successive companions have coordinates
\[
 a_{2r}:(i,p-r)\to(j,q-r),\qquad
 a_{2r+1}:(j,q-r-1)\to(i,p-r).
\]
Complete each maximal open companion chain separately modulo $2M$, identifying the coordinates of an old periodic endpoint modulo its old period. Retain the chain label, so arrows from different chains are not identified merely because they have the same endpoints. If both endpoints belong to finite intervals, any new arrow with both endpoints old would require, in the even case,
$u-v=p-q+kM$. The difference between the two sides with $k=0$ has absolute value at most $2L$; hence $k=0$. For odd arrows the analogous bound is $2L+1$. Before reduction modulo $M$, the even terms with old endpoints have
\[
 \max(p-l_i,q-l_j)\le r\le\min(p,q),
\]
and the odd terms have
\[
 \max(p-l_i,q-l_j-1)\le r\le\min(p,q-1).
\]
Each endpoint of the odd range equals the corresponding endpoint of the even range or that endpoint minus one. Thus the allowed even and odd indices interlace to form one consecutive interval of companion indices. Forward and inverse companion completeness puts every term of this interval in the original chain. Thus no new old-to-old arrow occurs. With one periodic endpoint, the finite endpoint determines the unique old occurrence, while divisibility of $M$ by the old period preserves the other phase. With two periodic endpoints the chain was already closed.

Transport the vertex fields along the completed cycles. For a completed interval, choose the transport at one new closing edge to be the inverse of the product of all other transports. Its full field transport is then the identity; for old cycles this holds after $M$ steps by the choice of $M$. Extend each arrow bimodule by alternating duality relative to its two endpoint fields, using the original nondegenerate pairings wherever prescribed. After $2M$ companions the endpoint field actions agree with the original ones. Choose the last new pairing to close the chain and retain the resulting invertible linear return map as part of the species transport $F$. No identity condition on this linear return map is needed.

On each open companion chain, alternate the forward and reverse designation. Any descent automorphism preserves its two ends and commutes with taking companions, so it preserves this designation; a shift along a finite chain cannot preserve its ends. On the old closed chains use the Dynkin cuts, including the central exceptional cut, just constructed. The double-dual identifications give a species transport $F$. Choose the newly added mesh pairings to be the transpose pairings. On an old open chain, the prescribed pairings determine these choices successively. In split one-dimensional arrow coordinates, if the terms at the consecutive meshes have nonzero coefficients $c_{2r}$ and $c_{2r+1}$, take the forward arrow to be $a_{2r}$, replace the reverse arrow $a_{2r+1}$ by $c_{2r}a_{2r+1}$, and set $F(a_{2r})=-c_{2r+1}c_{2r}^{-1}a_{2r+2}$. The transpose terms then have the original coefficients at both meshes. For general bimodules the same construction uses the two duality isomorphisms in place of the scalar ratio. An open chain imposes no return condition on these choices. On the closed components the preceding Dynkin-coordinate calculation supplies the compatible choice. This gives the required transpose relations and the exceptional square relations.

Finally quotient by all new vertex idempotents. Each relation at a new vertex vanishes, since its endpoint is new. At an old vertex whose predecessor was missing, the relation has a new translated endpoint and vanishes as well. Every other old relation loses exactly the terms with new intermediate vertices. It is therefore precisely its original partial mesh relation. The image of an ideal under a quotient map is generated by the images of its generators. Thus all relations in the vertex quotient are generated by these images of mesh relations, and the original and quotient algebras have identical presentations. An added cycle contains a new vertex because $M>L+1$.
\end{proof}

The construction also applies to the opposite category. Reverse arrows and translation, and replace each bimodule by its opposite with the interchanged endpoint actions. Pairings and the square-boundary equations are carried to the corresponding opposite equations.

\subsection{Hall polynomials in the standard case}

\begin{proof}[Proof of Theorem~\ref{thm:standard}]
Apply the preceding presentation and completion to the opposite of $\Lambda$, so that its covariant representations are the right defect modules used in Section~\ref{sec:defects}. A representation of the quotient is exactly a representation of the completion with zero spaces at all new vertices. Each added translation cycle contains such a vertex, and hence cannot support a nonzero local band. Thus the full cycles supporting bands are exactly the original periodic components.

Fix a conservative extension $E/k$ of $A$. Its degree is coprime to every residue-field degree $d_i$ from Section~\ref{sec:defects}. All vertex fields of the presentation and completion are copies of these fields. In the ordinary periodic components, Lemma~\ref{lem:residues} shows that no extra extension condition is required to preserve the indecomposables of the Dynkin species. Equation~\eqref{eq:ordinarytotal} gives their total masses in $\R$. On an exceptional component, Proposition~\ref{prop:squaretotal} gives the same conclusion; replacing its field-size variable by $Q^d$ preserves $\R$.

Proposition~\ref{prop:recovery} now gives all dimension, arrow-top and arrow-socle masses supported on the original vertices. On $\Lambda$ the arrow ideal is the radical, since the mesh grading has nilpotent positive part and semisimple degree zero. These are therefore the actual module top and socle. At vertex $i$, a defect dimension $\delta_i$ measured over $k$ is dimension $\delta_i/d_i$ over its residue field; a nonintegral quotient means that the corresponding mass is zero. We have obtained exactly the masses required by the defect criterion, on every conservative extension of $A$. Corollary~\ref{cor:criterion} proves the theorem.
\end{proof}

\subsection{Ray-standard forms}
Over an algebraically closed field, a basic algebra is \emph{distributive} when its lattice of two-sided ideals is distributive. Its ray category has primitive idempotents as objects and, in each corner, the orbits under the two corner unit groups as rays. For composable morphisms $a:x\to y$ and $b:y\to z$, the ray product is zero when $ba\in b\rad A(y,y)a$, and otherwise is the ray of $ba$. Here $ba$ denotes categorical composition. The linearization has the nonzero rays as a basis and extends this product bilinearly. An algebra is \emph{ray-standard} when it is isomorphic to this linearization. A category is \emph{mild} when every proper ideal quotient is locally representation-finite.

A contour is a pair of paths with the same nonzero ray. It is essential if the equality cannot be obtained by replacing equal-ray proper subpaths, and non-deep if the common ray is not annihilated by both endpoint corner radicals. The full subcategory on its vertices is its support category. The structure and disjointness theorems in \cite[Theorem 6]{Bongartz13} identify the essential non-deep support categories as pennies, dumb-bells and diamonds, with distinct supports disjoint. The penny path coordinates needed here are specified in Section~\ref{sec:boundary}. In the stem category, morphisms are the subbimodules generated by rays and composition is their product. The mixed forms of \cite[Theorem 10]{Bongartz13} use stem composition on selected penny supports and ray composition elsewhere.

\begin{lemma}\label{lem:ray-mesh}
A ray-standard representation-finite algebra over an algebraically closed field is mesh-standard.
\end{lemma}
\begin{proof}
Let $\Gamma$ be its Auslander--Reiten quiver and $\mathcal C=k(\Gamma)$ its mesh category. By \cite[Theorem 14]{Bongartz13}, the universal mesh category covers the category of indecomposables. This covering preserves radicals. Indeed, if a noninvertible map upstairs had invertible image, lift an inverse using the direct-sum Hom isomorphism in the definition of a covering. The component returning to the original object gives a left inverse. Local endomorphism rings then make the resulting nonzero idempotent on the other object the identity, so the original map would be invertible. The contradiction proves the assertion.

The finite category of indecomposables has nilpotent radical, say of exponent $N$. Every path of length $N$ upstairs maps to zero, hence is zero, since the covering is injective on each individual Hom space. Every path in $\mathcal C$ lifts, so its length-$N$ paths vanish too. As $\Gamma$ is finite, $\mathcal C$ is finite-dimensional; the universal mesh map to $\mathcal C$ is a covering between locally bounded categories.

Apply \cite[Theorem 15]{Bongartz13} first to the covering from the universal mesh category to the indecomposable category of the given algebra. The universal mesh category is therefore an Auslander category. Apply the same theorem to its covering of $\mathcal C$; hence $\mathcal C$ is also an Auslander category and is the indecomposable category of its algebra of projective points. These points are detected by the projective dimensions of simple contravariant functors: the radical inclusion at a projective point gives a resolution of length at most one, whereas at a nonprojective point the minimal mesh resolution has a nonzero term in degree two. The projective points are therefore exactly those of $\Gamma$. By \cite[Theorem 18(c),(d)]{Bongartz13}, their algebra is the ray-standard form, which is the given algebra.
\end{proof}



\begin{thebibliography}{99}

\bibitem{BTCB}
R. Bennett--Tennenhaus and W. Crawley-Boevey,
\emph{Semilinear clannish algebras},
arXiv:2204.12138v3, 2022.

\bibitem{Bongartz13}
K. Bongartz,
\emph{On representation-finite algebras and beyond},
arXiv:1301.4088v2.

\bibitem{Bongartz23}
K. Bongartz,
\emph{On representation-finite selfinjective algebras, coverings, multiplicative bases \dots},
arXiv:2305.17418v1.

\bibitem{Hubery}
A. Hubery,
\emph{Ringel--Hall algebras of cyclic quivers},
S\~ao Paulo J. Math. Sci. \textbf{4} (2010), 351--398.

\bibitem{Jiao}
P. Jiao,
\emph{Auslander's defect formula and a commutative triangle in an exact category},
Front. Math. China \textbf{15} (2020), 115--125.

\bibitem{Lang}
S. Lang,
\emph{Algebraic groups over finite fields},
Amer. J. Math. \textbf{78} (1956), 555--563.

\bibitem{Riedtmann}
C. Riedtmann,
\emph{Algebren, Darstellungsk\"ocher, \"Uberlagerungen und zur\"uck},
Comment. Math. Helv. \textbf{55} (1980), 199--224.

\bibitem{Ringel}
C. M. Ringel,
\emph{Hall polynomials for the representation-finite hereditary algebras},
Adv. Math. \textbf{84} (1990), 137--178.

\end{thebibliography}

\medskip
\begingroup\small
\noindent\textsc{Zhili College, Tsinghua University, Beijing 100084, China}\par
\noindent\textit{Email address:}\enspace\mbox{\href{mailto:tangqy24@mails.tsinghua.edu.cn}{\texttt{tangqy24@mails.tsinghua.edu.cn}}}\par
\endgroup
\end{document}
